\documentclass[11pt,a4paper]{article}
\usepackage[T1]{fontenc}
\usepackage[utf8]{inputenc}
\usepackage{lmodern}
\usepackage[margin=25mm]{geometry}
\usepackage{amsmath,amssymb,booktabs,longtable,array}
\usepackage{xcolor,graphicx,listings}
\usepackage{hyperref}
\hypersetup{colorlinks=true,linkcolor=blue!45!black,urlcolor=blue!45!black,
  citecolor=blue!45!black,pdftitle={Kronos Financial Forecasting Lab: Technical Report},
  pdfauthor={Kronos Financial Forecasting Lab}}
\lstset{basicstyle=\ttfamily\footnotesize,breaklines=true,columns=fullflexible,
  keepspaces=true,frame=single,rulecolor=\color{black!20},showstringspaces=false}
\setlength{\emergencystretch}{2em}
\setlength{\parskip}{4pt}
\setlength{\parindent}{0pt}
\renewcommand{\arraystretch}{1.15}
% Generated by generate_tables.py from evidence/run-summary.json.
\newcommand{\BaseMAE}{185.40}
\newcommand{\BaselineMAE}{290.77}
\newcommand{\MAEReductionPercent}{36.24\%}
\newcommand{\RawCoverage}{65.625\%}
\newcommand{\CalCoverage}{96.875\%}
\newcommand{\RawWidth}{418.72}
\newcommand{\CalWidth}{2953.80}
\newcommand{\WidthRatio}{7.05}
\newcommand{\RawPinball}{65.52}
\newcommand{\CalPinball}{131.18}

\newcommand{\code}[1]{\texttt{\detokenize{#1}}}
\newcommand{\sourcecommit}{693bcc4331a2df80402496492ccfa1a237f90c7a}
\title{Kronos Financial Forecasting Lab\\[5pt]
  \large Implementation, chronological evaluation, and adaptation evidence}
\author{Technical report for \texttt{tuanngocfun/agent}\\
  \small Branch: \texttt{kronos-evaluation-lab}}
\date{25 October 2026 (Asia/Ho\_Chi\_Minh)}

\begin{document}
\maketitle
\begin{abstract}
The lab implements real Kronos inference with retained Monte Carlo paths,
chronological evaluation against persistence, empirical interval calibration,
two adaptation modes, a cost-aware paper simulator, and shared CLI, HTTP, and MCP
interfaces. This report documents the implemented methods and the recorded
24--25 October 2026 integration evidence. On four held-out BTCUSDT hourly windows
(32 origin--horizon observations), Kronos-base close MAE was \BaseMAE{} USDT,
compared with \BaselineMAE{} for persistence. Raw q10--q90 coverage was
\RawCoverage{}; earlier-label interval expansion raised coverage to
\CalCoverage{} while widening the mean interval from \RawWidth{} to
\CalWidth{} USDT and increasing pinball loss. A separate Kronos-small ablation
executed two optimizer updates per adapted stage. These measurements establish
that the workflows execute and preserve their intended data boundaries. Their
size, temporal dependence, and unaudited pretraining provenance do not support a
general forecast-skill, adaptation-gain, or profitability claim.
\end{abstract}

\newpage
\tableofcontents
\newpage

\section{Scope and evidence}
The deliverable is an executable research lab for the financial project in the
supplied research plans, rather than a completed broad-market research study.
The repository contains searchable versions of those plans under
\code{docs/plans/}; \code{docs/IMPLEMENTATION.md} maps their requirements to code
and states the corrections made during implementation~\cite{implementation}.
The implemented scope includes public/local OHLCVA ingestion, probabilistic
forecasting, chronological calibration and testing, checkpointed adaptation,
model comparison, local paper simulation, and agent access.

The quantitative source for this report is the committed
\code{evidence/run-summary.json}, recorded at 2026-10-24 17:20:28 UTC
(25 October, 00:20:28 in Ho Chi Minh City)~\cite{evidence}.
Implementation behavior is described at commit \texttt{\sourcecommit};
the availability fix and its four additional tests are documented separately in
\code{VALIDATION.md}~\cite{validation}. Tables and headline numerical values are
generated from the evidence file by \code{generate_tables.py}, and an accompanying
manifest records its SHA-256. Full trajectories, local data, and training
checkpoints remain local; they are not required to regenerate this document's
aggregate tables. Re-running the models requires those inputs or a fresh data
collection and can produce different measurements.

The dashboard screenshot supplied with the report request shows an interactive
configuration of 128 history bars, 12 future bars, and 20 paths. The numerical
evaluation below uses the recorded 64/8/8 experiment. A displayed live forecast
does not provide realized future labels or substitute for that evaluation.

\section{System design}
The same forecasting engine serves a Python CLI, an authenticated local HTTP API,
and a stdio Model Context Protocol (MCP) server. The main boundaries are shown in
Figure~\ref{fig:architecture}. Heavy jobs run through a bounded shared executor;
inference is serialized to limit resource contention and shared random-state
effects. The HTTP process stores jobs and results in bounded memory, so users
export results before restarting it.

\begin{figure}[htbp]
\centering
\fbox{\parbox{0.91\linewidth}{\centering
Public Binance / local CSV / optional Yahoo and Vnstock adapters\\[4pt]
$\downarrow$\\[4pt]
UTC validation and provenance $\longrightarrow$ history-only normalization\\[4pt]
$\downarrow$\\[4pt]
Pinned tokenizer + Kronos $\longrightarrow$ retained $S\times H\times6$ paths\\[4pt]
$\downarrow$\\[4pt]
CLI / authenticated HTTP jobs / stdio MCP\\[4pt]
$\downarrow$\\[4pt]
Dashboard and JSON export; chronological scoring and paper simulation}}
\caption{Implemented data and execution boundaries. Observed future prices enter
the scorer after a forecast has been produced. Known future timestamps may be
passed to the forecasting engine.}
\label{fig:architecture}
\end{figure}

\begin{table}[htbp]
\centering\small
\begin{tabular}{@{}>{\raggedright\arraybackslash}p{0.26\linewidth}>{\raggedright\arraybackslash}p{0.66\linewidth}@{}}
\toprule
Module & Responsibility \\
\midrule
\code{data.py} & Provider adapters, strict frame validation, amount provenance. \\
\code{forecast.py} & Checkpoint identities, normalization, autoregressive sampling, path summaries. \\
\code{evaluation.py} & Matching origins, disjoint labels, frozen calibration, metrics. \\
\code{training.py}, \code{ablation.py} & Causal windows, stage training, checkpoint selection, paired variants. \\
\code{comparison.py}, \code{paper.py} & Chronos-2 comparison and cost-aware LONG/FLAT simulation. \\
\code{service.py}, \code{api.py}, \code{mcp_server.py} & Bounded jobs and shared access to model evidence. \\
\bottomrule
\end{tabular}
\caption{Source modules under \texttt{kronos\_lab/} at the implementation commit.}
\end{table}

\section{Data contract and provenance}
Each candle has a UTC opening timestamp and six ordered features:
\[
  x_t=(O_t,H_t,L_t,C_t,V_t,A_t),
\]
where $A_t$ is quote amount. The recorded Binance integration uses actual quote
volume and closed candles. When another supported input omits amount, the adapter
labels its proxy, $A_t=V_t(O_t+H_t+L_t+C_t)/4$. Missing or invalid prices and
volumes, duplicate timestamps, and inconsistent observed candles are rejected.
Generated candles are treated separately: their positivity and OHLC consistency
violations are counted and exposed, rather than silently repaired.

The saved input consists of 499 closed BTCUSDT hourly candles, with opening
timestamps from 3 October 2026 22:00 UTC through 24 October 2026 16:00 UTC. All
tables use UTC for data boundaries. The CSV identity is:
\begin{quote}\small\ttfamily
45d795760f542ce7bac9af39331b2c0b91f417c16c0d159741af2faf1d022d2d
\end{quote}
This identifies the saved input, rather than promising that a later public API
request will retrieve an identical file.

Future timestamps must be unique, increasing, and later than the history. A
constant observed cadence can be extended for continuously traded data; it is
not an exchange-session calendar. Provider-backed equity forecasts require
explicit future sessions, and a regular-looking imported CSV does not establish
weekend or holiday trading. Yahoo uses unadjusted prices and records that policy.
Yahoo and Vnstock integrations were exercised through mocked adapter tests;
their live providers were not part of the recorded experiment.

\section{Forecast construction}
\subsection{History-only scaling and discrete generation}
Let $L$ denote the lookback and $H$ the horizon. For each feature $j$, the engine
fits population statistics on the $L$ observed history rows only:
\begin{align}
\mu_j&=\frac{1}{L}\sum_{t=1}^{L}x_{t,j}, &
\sigma_j&=\sqrt{\frac{1}{L}\sum_{t=1}^{L}(x_{t,j}-\mu_j)^2},\\
z_{t,j}&=\operatorname{clip}\!\left(\frac{x_{t,j}-\mu_j}{\sigma_j+10^{-5}},-5,5\right).
\end{align}
The official Kronos design uses a tokenizer for hierarchical discrete K-line
tokens and an autoregressive transformer for their evolution~\cite{kronos}.
The lab pins the upstream code and model revisions (Appendix~\ref{app:identities})
and requires a clean upstream checkout. It encodes the observed history once,
replicates its token sequence across $S$ samples, and samples each next token
half from the predictor's two dependency-aware heads. Calendar features are
minute, hour, weekday, day, and month.

Logits are divided by temperature before optional top-$k$ and nucleus top-$p$
filtering. Decoded future paths are transformed back using the same historical
statistics. The retained tensor is
\[
  \widehat X\in\mathbb R^{S\times H\times6}.
\]
The implementation enforces $L+H\leq512$ and $1\leq S\leq256$; it does not
discard paths by averaging them before uncertainty analysis. The recorded
Kronos-base run used $L=64$, $H=8$, $S=8$, temperature $0.6$, top-$p=0.9$,
top-$k=0$, and seed 42 on an A100 MIG 20 GB device, with Python 3.12.3 and
PyTorch 2.6.0+cu124. Its returned shape was $(8,8,6)$, with zero diagnosed invalid
generated candles. This one diagnostic outcome does not guarantee that future
generated candles will always be valid.

\subsection{Point summaries and sample votes}
For close-price paths $\widehat C_{s,h}$ and last observed close $c_i$ at origin
$i$, the point forecast and sample vote are
\begin{align}
\widehat c_{i,h}&=\frac{1}{S}\sum_{s=1}^{S}\widehat C_{i,s,h}, &
\widehat p_{i,h}&=\frac{1}{S}\sum_{s=1}^{S}\mathbf 1\{\widehat C_{i,s,h}>c_i\}.
\end{align}
The API exposes the terminal vote as \code{p_up} and q10/q50/q90 as empirical
close quantiles. The expected terminal return is
$100(S^{-1}\sum_s\widehat C_{i,s,H}/c_i-1)$ percent. Sample votes describe the
generated path cloud; they are not calibrated correctness probabilities.
With eight paths their frequency resolution is 0.125. A smooth density display
does not increase the number of model draws or establish calibrated uncertainty.
No universal volatility multiplier or protocol-derived confidence score is added.

\section{Chronological evaluation and interval calibration}
\subsection{Origins and information boundaries}
At origin index $i$ (the first unknown row), the forecaster receives rows
$[i-L,i)$ and the future timestamps $[i,i+H)$; future OHLCVA values remain in the
scorer. Persistence predicts $c_i$ at every horizon on exactly the same origins.
Calibration labels precede test labels. Overlapping cross-split labels are
purged, and the recorded run reports disjoint labels and no within-split label
overlap. Historical contexts may overlap, so label disjointness does not imply
statistical independence.

\begin{table}[htbp]
\centering\footnotesize
% Generated by generate_tables.py from evidence/run-summary.json.
\begin{tabular}{lrrll}
\toprule
Split & First origin & Last origin & Label start & Label end \\
\midrule
Calibration & 432 & 456 & 2026-10-21 22:00 & 2026-10-23 05:00 \\
Test & 464 & 488 & 2026-10-23 06:00 & 2026-10-24 13:00 \\
\bottomrule
\end{tabular}

\caption{Main Kronos-base evaluation. Origin indices are zero-based; all label
times are UTC. Each split has four eight-bar windows. Test origins are
464, 472, 480, and 488.}
\label{tab:splits}
\end{table}

\subsection{Empirical additive expansion}
For calibration window $i$ and horizon $h$, let $\ell_{i,h}$ and $u_{i,h}$ be
the raw q10 and q90 limits, and $y_{i,h}$ the realized close. The nonnegative
score is
\[
 a_{i,h}=\max\{\ell_{i,h}-y_{i,h},\;y_{i,h}-u_{i,h},\;0\}.
\]
For $n$ calibration windows, the selected order-statistic rank is
$k=\min\{n,\lceil(n+1)0.8\rceil\}$. The $k$th sorted score at each horizon
defines $\delta_h$. It is frozen before testing and applied as
\[
 [\ell^{\mathrm{cal}}_{i,h},u^{\mathrm{cal}}_{i,h}]
   =[\ell_{i,h}-\delta_h,\;u_{i,h}+\delta_h].
\]
The point mean, median, and votes are unchanged. Here $n=4$ and $k=4$, so each
expansion is the largest observed calibration miss for its horizon. The rank
does not exceed $n$ in this run; the code reports a cap when it does. This is an
empirical chronological use of split-conformal-style scores, with no claimed
exchangeability-based coverage guarantee for serially dependent market data.

\begin{table}[htbp]
\centering\small
% Generated by generate_tables.py from evidence/run-summary.json.
\begin{tabular}{rr}
\toprule
Horizon & Expansion \\
\midrule
1 & 210.60 \\
2 & 574.41 \\
3 & 627.20 \\
4 & 1279.27 \\
5 & 1821.58 \\
6 & 1797.17 \\
7 & 1871.65 \\
8 & 1958.43 \\
\bottomrule
\end{tabular}

\caption{Frozen Kronos-base interval expansions, in USDT, fitted on the earlier
four calibration windows.}
\end{table}

\subsection{Metric definitions}
Let $N=WH$ for $W$ test origins, $e_{i,h}=\widehat c_{i,h}-y_{i,h}$, and
$r_{i,h}=100(y_{i,h}/c_i-1)$. The reported close errors are
\[
 \mathrm{MAE}=N^{-1}\sum_{i,h}|e_{i,h}|,\qquad
 \mathrm{RMSE}=\sqrt{N^{-1}\sum_{i,h}e_{i,h}^2}.
\]
Return MAE uses $|\widehat r_{i,h}-r_{i,h}|$ and is measured in percentage points.
Directional accuracy compares the signs of price changes from the observed
anchor, including a distinct zero-change class. Persistence therefore has zero
directional accuracy here: its prediction is flat and all scored actual changes
are nonzero. This is not a binary up/down accuracy of zero.

Brier score is $N^{-1}\sum_{i,h}(\widehat p_{i,h}-\mathbf1\{y_{i,h}>c_i\})^2$.
Persistence is a point forecast without a probability distribution; a separate
constant probability 0.5 benchmark has Brier score 0.25. Coverage is the fraction
of observed closes inside the inclusive q10--q90 limits; width is the mean
difference between their bounds. For residual $u=y-q$, pinball loss is
$\rho_\tau(u)=\max\{\tau u,(\tau-1)u\}$. The reported mean pinball loss
averages the losses for $\tau=0.1,0.5,0.9$, using the unchanged median for q50.
Terminal metrics score only $h=H$ (four observations here), whereas the main
tables pool all 32 origin--horizon observations.

\section{Recorded results}
\subsection{Kronos-base point forecast and uncertainty}
\begin{table}[htbp]
\centering\small
% Generated by generate_tables.py from evidence/run-summary.json.
\begin{tabular}{lrrrrr}
\toprule
Method & MAE & RMSE & Return MAE & Direction acc. & Brier \\
\midrule
Persistence & 290.77 & 365.17 & 0.36 & 0.00\% & -- \\
Kronos-base & 185.40 & 259.46 & 0.23 & 81.25\% & 0.16309 \\
\bottomrule
\end{tabular}

\caption{Kronos-base test point metrics against matching persistence.
MAE/RMSE are USDT; return MAE is percentage points; direction is percent;
Brier is unitless. Expansion leaves these model quantities unchanged.}
\label{tab:point}
\end{table}
Kronos-base MAE is \MAEReductionPercent{} lower than persistence on these
four windows. Its terminal directional accuracy is 75\% (three of four), and
terminal Brier score is 0.13672. These are descriptive comparisons on one short
BTCUSDT segment. Thirty-two correlated horizon observations are not 32
independent experiments, and there is no confidence interval or significance
claim attached to the observed MAE difference.

\begin{table}[htbp]
\centering\small
% Generated by generate_tables.py from evidence/run-summary.json.
\begin{tabular}{lrrr}
\toprule
Method & Coverage & Width & Pinball \\
\midrule
Raw & 65.625\% & 418.72 & 65.52 \\
Calibrated & 96.875\% & 2953.80 & 131.18 \\
\bottomrule
\end{tabular}

\caption{Kronos-base q10--q90 test intervals. Coverage is percent; width and
mean pinball loss are in USDT. The label ``80\%'' denotes quantile levels,
not verified future coverage.}
\label{tab:interval}
\end{table}
Raw intervals cover 21 of 32 closes; expanded intervals cover 31 of 32. The
mean width increases by a factor of \WidthRatio{}, and mean pinball loss
increases from \RawPinball{} to \CalPinball{}. The earlier calibration misses
therefore lead to very broad test intervals. Higher coverage alone does not
establish a better predictive distribution, and these results provide no
evidence that one expansion will transfer to other market regimes.

\subsection{Kronos-small adaptation ablation}
The adaptation experiment starts from Kronos-small, not Kronos-base. It uses
three variants with identical initial model/tokenizer identities, scored test
origins, and sampling settings: unchanged zero-shot; predictor-only adaptation
with a frozen original tokenizer; and tokenizer adaptation followed by predictor
adaptation with the selected tokenizer frozen. The recorded trained stages each
take two optimizer updates, in one epoch, with batch size two and seed 42.
Tokenizer/predictor learning rates are $2\cdot10^{-4}$ and $10^{-6}$, and
gradient clipping is 2.0 in this bounded run.

The four half-open partitions have exclusive upper boundaries 12 October
06:00, 16 October 10:00, 20 October 14:00, and 24 October 17:00 UTC. These
separate optimization, checkpoint-selection validation, interval calibration,
and final test. The ablation's scored calibration labels are 19 October 06:00
through 20 October 13:00 UTC; they differ from the main base-model calibration
in Table~\ref{tab:splits}. Its test labels match the main experiment. Training
dataset window starts and scored evaluation origins are separately constructed;
partition boundaries alone do not determine the exact scored label range.

\begin{table}[htbp]
\centering\small
% Generated by generate_tables.py from evidence/run-summary.json.
\begin{tabular}{llrrrr}
\toprule
Variant & Steps & MAE & RMSE & Raw cov. & Cal. cov. \\
\midrule
zero-shot & 0 & 214.55 & 282.10 & 59.375\% & 93.750\% \\
predictor-only & 2 P & 214.55 & 282.10 & 59.375\% & 93.750\% \\
two-stage & 2 T + 2 P & 213.49 & 281.16 & 59.375\% & 90.625\% \\
\bottomrule
\end{tabular}

\caption{Separate Kronos-small ablation. P/T mean predictor/tokenizer optimizer
updates. MAE/RMSE are USDT; raw and expanded coverage are percent. All rows use
four held-out test windows.}
\label{tab:ablation}
\end{table}
Predictor-only and zero-shot metrics are identical at the reported precision
(and in the saved aggregate values) despite a changed saved predictor hash.
Two-stage MAE is slightly lower in this run, but two updates and four test
windows cannot identify a stable adaptation gain. Changed checkpoints establish
that training executed; they do not establish a meaningful forecast improvement.
The base and small tables are not a controlled comparison of model capacity:
their calibration windows differ and this experiment was designed for integration
validation rather than capacity selection.

\subsection{Chronos-2 and paper simulation}
The comparator loads \code{amazon/chronos-2} on CPU with close as the target and
past volume as a covariate. It uses the same history and future timestamps as
the compared Kronos forecast. The recorded terminal directions disagree. This
is one paired live forecast, not a held-out Chronos-2 skill benchmark, and
agreement or disagreement cannot by itself establish which forecast is accurate.

The paper simulator chooses LONG only when the frozen terminal mean forecast
exceeds the last observed close; otherwise it stays FLAT. It enters at the next
candle's open, exits at the final forecast candle's close, and rejects overlapping
test windows. For cash $E$, fee fraction $f$, adverse slippage $s$, entry open
$O$, and exit close $C$, a LONG roundtrip uses
\[
 Q=\frac{E}{O(1+s)(1+f)},\qquad E'=Q\,C(1-s)(1-f).
\]
Fees apply to each actual fill notional. Drawdown is sampled at the entry mark,
held closes, and exit; it does not measure intrabar lows. The recorded run uses
10 bps fees and 5 bps slippage per side and stays FLAT on all four test windows:
zero trades, zero turnover, zero costs, and zero return. Consequently it does
not exercise active-trade behavior on real saved test windows or demonstrate a
profitable strategy. Active-trade accounting is covered by dedicated tests.

\section{Adaptation implementation and leakage controls}
\subsection{Window construction and objectives}
Each training window and its labels are fully contained in one partition; the
dataset does not borrow context from the preceding partition. It fits each
feature's mean and population standard deviation on the window's context only,
then transforms both context and target using those statistics, epsilon
$10^{-5}$, and clipping at five. This prevents target values from changing the
normalization of observed history.

Tokenizer adaptation minimizes the exact implemented objective
\[
 \mathcal L_{\mathrm{tok}}=\tfrac12\bigl[
  \mathrm{MSE}(\widehat z^{\mathrm{coarse}},z)+
  \mathrm{MSE}(\widehat z^{\mathrm{full}},z)+\mathcal L_{\mathrm{quant}}\bigr],
\]
where the quantization term comes from the pinned upstream tokenizer. Predictor
adaptation freezes the tokenizer and applies the average of its two head
cross-entropies, $(\mathrm{CE}_{s1}+\mathrm{CE}_{s2})/2$, to
future positions only, starting at the context boundary. Teacher forcing uses
the shifted target token sequence. Loss on already observed context positions
is excluded. Both stages save their best validation checkpoint and restore it
before the next stage or evaluation; calibration and test labels do not select
checkpoints~\cite{training}.

\subsection{Causal validation and checkpoint identity}
The pinned upstream cross-attention selects causal behavior through its module
training flag. A plain evaluation-mode teacher-forced validation pass could
therefore expose later sibling-token information. The lab's
\code{causal_loss_mode} preserves causal attention during validation while
disabling attention dropout and restoring prior module state afterward. This
is a concrete guard in the implemented validation path, not a claim that every
possible dependency in pretrained representations has been causally audited.

Hub references resolve to immutable model SHAs; local checkpoint identities
hash the configuration and weight files. Training manifests preserve initial
and selected checkpoint identities and the partition policies. The zero-shot,
predictor-only, and two-stage evaluations use the same scored origins. Tests
mutate held-out data to check that checkpoint selection does not change. The
separate two-worker CPU distributed smoke uses small randomly initialized
official architectures and one update per stage: it verifies distributed
execution and saving, not pretrained multi-GPU training or speedup.

\section{Interfaces, operation, and verification}
The MCP server exposes these tools:
\begin{lstlisting}
forecast_kline, forecast_signal, forecast_compare
forecast_batch, job_status
\end{lstlisting}
Heavy
work runs outside its protocol event loop; batch outcomes isolate invalid
symbols. MCP transports model evidence and does not assign a statistical
confidence level. A separate news corpus, retrieval pipeline, or multimodal
reasoning model is outside this implementation.

HTTP jobs use \code{POST /api/jobs}, \code{POST /api/batch}, and
\code{GET /api/jobs/{job_id}}, with bearer-token authentication for model
access. The token is read locally and is never part of this document. Server
errors omit raw traces and secret values. The dashboard supports forecast paths,
raw quantiles, empirical 3D density, held-out evaluation, Chronos-2 comparison,
CSV/saved-result import, and JSON export. Its raw q10--q90 display must be
interpreted separately from held-out calibrated coverage.

The persistent \code{scripts/dashboard.sh} start/status/stop helper anchors
paths to the checkout, waits for a successful health check, and tracks the
owned process under \code{.cache/dashboard/}. It survives terminal closure
but is not configured to restart at machine reboot. The availability fix was
verified with HTTP 200 responses and a newly completed authenticated Kronos-base
GPU job, and added four lifecycle tests. It does not imply that the service will
remain available after process termination or host restart.

\begin{table}[htbp]
\centering\small
\begin{tabular}{@{}p{0.30\linewidth}p{0.62\linewidth}@{}}
\toprule
Check & Recorded verification and scope \\
\midrule
Project tests & 155 in the original integration record; 159 after the four lifecycle tests, also re-run for this report (159 passed in 14.26 s; one nonfatal timedelta deprecation warning). \\
Pinned upstream & Four official CPU regression tests with real small-model/tokenizer weights. \\
Real inference & Kronos-base GPU paths, chronological scoring, pretrained small-model adaptation, CPU Chronos-2. \\
Distributed execution & Two CPU workers, random-test initialization, one update per stage; no multi-GPU run. \\
MCP & Official SDK memory/stdio tests and real pretrained signal/batch success/error isolation. \\
Browser & Desktop/mobile, saved forecast/evaluation/comparison, 3D view, export, authenticated job, stale-session-state regression. \\
Paper accounting & Dedicated active-trade cost/slippage tests; recorded real windows all stayed FLAT. \\
\bottomrule
\end{tabular}
\caption{Validation evidence from \texttt{VALIDATION.md} and the run summary.
These checks verify software behavior within their tested scope.}
\end{table}

\section{Reproduction}
\subsection{Rebuild this report}
The LaTeX source, generated tables, generator, source-hash manifest, and compiled
PDF are stored together under \code{docs/technical-report/}. The table generator
uses only the Python standard library and the committed aggregate evidence.
From the repository root:
\begin{lstlisting}
python3 docs/technical-report/generate_tables.py
python3 docs/technical-report/generate_tables.py --check
make -C docs/technical-report
\end{lstlisting}
The Makefile supports Tectonic or pdfLaTeX. The published build uses Tectonic
0.17.0. A first Tectonic build may download its TeX bundle; an offline rebuild
needs a populated cache or a complete TeX installation. Byte-for-byte PDF
identity is not promised across compiler/bundle versions. Table and manifest
generation is deterministic. The report README records the build and inspection
commands and distinguishes document checks from previously recorded model runs.

\subsection{Run new forecasts and evaluations}
The application environment is separate from the report toolchain. The setup
script creates the Python environment and pinned upstream checkout; model
weights may be downloaded on first use. \code{requirements.lock} records the
tested Linux environment. From the project root, representative commands are:
\begin{lstlisting}
bash scripts/setup.sh
.venv/bin/python -m kronos_lab download \
  --symbol BTCUSDT --interval 1h --limit 1000
.venv/bin/python -m kronos_lab forecast \
  --csv data/BTCUSDT_1h.csv --device cuda:1 \
  --lookback 64 --horizon 8 --sample-count 8 \
  --output output/new-forecast.json
.venv/bin/python -m kronos_lab evaluate \
  --csv data/BTCUSDT_1h.csv --device cuda:1 \
  --lookback 64 --horizon 8 --sample-count 8 \
  --max-windows 8 --output output/new-evaluation.json
.venv/bin/python -m pytest tests -q
\end{lstlisting}
Choose a device that is actually available. A new download changes the time
range and resulting split; these commands execute the workflow rather than
recovering the exact saved October labels. Exact replay additionally requires
the saved CSV hash, immutable checkpoint revisions, sampling settings, evaluation
origin policy, and compatible numerical environment. GPU numerical behavior
can vary even with the same seed. Research adaptation should use a new output
directory and declared boundaries; the recorded bounded configuration is
available locally under \code{output/real-run/ablation-config.json}.

\section{Limits and research work remaining}
The supported conclusion is that an integrated lab runs real pretrained models,
preserves trajectories and data boundaries, and reports inspectable evidence.
The main unresolved scientific question is whether any improvement persists
across origins, seeds, assets, market regimes, and realistic trading costs.
Four calibration windows make interval expansion unstable; four test windows
and two training updates cannot settle that question. Pretraining data coverage
was not independently audited, so an immutable checkpoint does not establish
an uncontaminated historical benchmark.

A subsequent study needs a declared multi-asset/time-range benchmark, sufficient
non-overlapping origins, a fixed checkpoint-selection/calibration/test protocol,
longer adaptation with validation-based stopping, multiple seeds, and uncertainty
estimates that respect temporal dependence. It should compare skill and interval
sharpness jointly, retain negative results, and test active trading with costs
over enough trades to assess behavior. A rolling or regime-specific calibration
study would address the extreme interval widening seen here. Such experiments
are not represented as already completed in this report. Live Yahoo/Vnstock,
pretrained multi-GPU training, broad equities, cross-domain/biomedical transfer,
and news/multimodal retrieval likewise remain unrun or separate projects.

\appendix
\section{Immutable identities}\label{app:identities}
The following identifiers come from the committed evidence~\cite{evidence}.
They identify source code and weights; they do not audit model pretraining data.

\begin{description}
\item[Lab implementation commit]\hfill\\
\code{693bcc4331a2df80402496492ccfa1a237f90c7a}
\item[Official Kronos code]\hfill\\
\code{67b630e67f6a18c9e9be918d9b4337c960db1e9a}
\item[\texttt{NeoQuasar/Kronos-base}]\hfill\\
\code{2b554741eca47781b64468546e77fef3e85130e6}
\item[\texttt{NeoQuasar/Kronos-small}]\hfill\\
\code{901c26c1332695a2a8f243eb2f37243a37bea320}
\item[\texttt{NeoQuasar/Kronos-Tokenizer-base}]\hfill\\
\code{0e0117387f39004a9016484a186a908917e22426}
\item[\texttt{amazon/chronos-2}]\hfill\\
\code{29ec3766d36d6f73f0696f85560a422f50e8498c}
\end{description}

The saved predictor-only model content hash is:
\begin{quote}\small\ttfamily
ac5f5ded4e2f9e2c7772c5ada72c20b67b46b17e8ee54e72505d73c261896d50
\end{quote}
\begin{samepage}
The saved two-stage tokenizer and model content hashes are, respectively:
\begin{quote}\small\ttfamily
f504eaf5cc3bf76208f283ec26427193805024e3bca35c019355421e9d1211c2\\
eb2b7aba4bc6526b3cf2f16a0a202eaa6330554b15f8cc136c54b2a4ea7aa0a3
\end{quote}
\end{samepage}
The manifest beside this report records the evidence-file hash and source
identities in machine-readable form. Raw candles, private tokens, model weights,
checkpoints, and dashboard screenshots are not embedded in the PDF.

\begin{thebibliography}{9}
\raggedright
\bibitem{evidence}
Kronos Financial Forecasting Lab. \emph{Recorded integration evidence},
24 October 2026, \code{evidence/run-summary.json} at lab implementation commit.
\href{https://github.com/tuanngocfun/agent/blob/693bcc4331a2df80402496492ccfa1a237f90c7a/evidence/run-summary.json}{Versioned aggregate evidence}.

\bibitem{implementation}
Kronos Financial Forecasting Lab. \emph{Plan-to-code mapping},
\code{docs/IMPLEMENTATION.md}; supplied plan text in \code{docs/plans/}.
\href{https://github.com/tuanngocfun/agent/blob/693bcc4331a2df80402496492ccfa1a237f90c7a/docs/IMPLEMENTATION.md}{Versioned implementation mapping}.

\bibitem{validation}
Kronos Financial Forecasting Lab. \emph{Validation and availability fix},
24--25 October 2026, \code{VALIDATION.md}.
\href{https://github.com/tuanngocfun/agent/blob/693bcc4331a2df80402496492ccfa1a237f90c7a/VALIDATION.md}{Versioned validation record}.

\bibitem{kronos}
Official Kronos project. \emph{Kronos: A Foundation Model for the Language of
Financial Markets}, README and source at audited revision
\code{67b630e67f6a18c9e9be918d9b4337c960db1e9a}.
\href{https://github.com/shiyu-coder/Kronos/tree/67b630e67f6a18c9e9be918d9b4337c960db1e9a}{Pinned upstream code}.

\bibitem{training}
Kronos Financial Forecasting Lab. \emph{Training, forecasting, and evaluation
implementation}, \code{kronos_lab/training.py}, \code{forecast.py}, and
\code{evaluation.py} at the lab implementation commit.
\href{https://github.com/tuanngocfun/agent/tree/693bcc4331a2df80402496492ccfa1a237f90c7a/kronos_lab}{Versioned application source}.
\end{thebibliography}
\end{document}
